% GENERATED FILE. Edit canonical lessons, then run npm run generate:books.
% canonical-source-hash: fnv1a64:e7ef236c1f30d923
% canonical-lessons: JA-C133-eki, JA-W133-ki, JA-C133-kinou, JA-C133-ike, JA-W133-ke, JA-C133-kesa, JA-C133-inu, JA-W133-nu, JA-C133-nuno, JA-C133-heya, JA-W133-he, JA-R133-last-four

\chapter{The Last Four Hiragana: Ki, Ke, Nu and He}
\label{ch:ja-last-four-hiragana-ki-ke-nu-he}
\hlchaptermodality{133}

\begin{chapteropening}
\textbf{By the end of this chapter:} \emph{I can write ki, ke, nu and he, the last four basic hiragana, and read words that use them, such as eki, inu and heya.}

Read eki, kinou, ike, kesa, inu, nuno and heya, write kinou and kesa from memory and say which is earlier, write heya with he at the front, and write the four new signs once each.
\end{chapteropening}

\section[\texorpdfstring{\ja{えき} (eki)}{えき (eki)}]{\ja{えき} (eki) — a station}

\label{lesson:JA-C133-eki}



\begin{quote}
Before the new word: write \textbf{\ja{え}} and say it. Then write \textbf{\ja{れ}} — \textbf{R2}, five lessons back.

\begin{itemize}
\raggedright
  \item \emph{Recall:} say \emph{a greengrocer} — \textbf{R4}, eighty lessons back
\end{itemize}
\end{quote}

\subsection*{You'll want to know: \ja{えき}}
\textbf{\ja{えき}} — \emph{eki} — \textquotedblleft{}a station\textquotedblright{}, the place where trains stop.

Two beats: \emph{e–ki}. The first sign is the \textbf{\ja{え}} you already write. The second one is new. It is one of only four hiragana signs the book has not yet written, and the next lesson writes it.

\subsection*{Your turn}
\begin{itemize}
\raggedright
  \item \emph{Say it:} \emph{eki}
  \item \emph{Say it:} \emph{eki}, clapping two beats
  \item \emph{Recall:} point to the sign in \textbf{\ja{えき}} you can already write, and the one you cannot
\end{itemize}

\begin{tcolorbox}[breakable,colback=teal!4,colframe=teal!35!black,title={Before you move on}]
What is \textquotedblleft{}a station\textquotedblright{}? (\textbf{\ja{えき}}, \emph{eki}.) How many beats? (\textbf{Two}.)
\end{tcolorbox}
\section[\texorpdfstring{\ja{き} (ki)}{き (ki)}]{\ja{き} — \ja{さ}, with a second bar}

\label{lesson:JA-W133-ki}



\begin{quote}
Four recalls, then the sign.

\begin{itemize}
\raggedright
  \item \emph{Recall:} say \emph{a station} — \textbf{R1}, one lesson back
  \item \emph{Recall:} say \emph{that}, a thing near the listener — \textbf{R2}, five lessons back
  \item \emph{Recall:} write \textbf{\ja{さ}}: the bar, the diagonal with its hook, the curve
  \item \emph{Recall:} say \emph{a barber} — \textbf{R4}, eighty lessons back
\end{itemize}
\end{quote}

\subsection*{Script — the shape}
\hlblockfigure{\detokenize{figures/JA-W133-ki-filmstrip.pdf}}{How \ja{き} is written, stroke by stroke}

\begin{quote}
\textbf{\ja{き}} — \emph{ki}
\end{quote}

Four strokes, three lifts:

\begin{enumerate}
\raggedright
  \item the upper \textbf{bar}, left to right
  \item the lower \textbf{bar}, left to right
  \item the long \textbf{diagonal} from the top, down to the right through both bars; at its foot, turn sharply and \textbf{hook} back to the left
  \item a separate \textbf{curve}: start at the left, drop down, and finish along the base to the right
\end{enumerate}

The diagonal, the hook and the curve are the ones your hand already makes in \textbf{\ja{さ}}. Only the top differs:

\noindent
\begin{tabularx}{\linewidth}{@{}>{\raggedright\arraybackslash}X>{\raggedright\arraybackslash}X@{}}
\toprule
\textbf{you write} & \textbf{bars at the top} \\
\midrule
\textbf{\ja{さ}}, \emph{sa} & one \\
\textbf{\ja{き}}, \emph{ki} & two \\
\bottomrule
\end{tabularx}

With \textbf{\ja{え}} in front of it, you can now write the word from the last lesson: \textbf{\ja{えき}}, \emph{eki}, a station.

\subsection*{Your turn}
\begin{enumerate}
\raggedright
  \item Trace \textbf{\ja{き}} saying \textquotedblleft{}bar, bar, diagonal and hook, curve\textquotedblright{}.
  \item Write \textbf{\ja{さ}} and \textbf{\ja{き}} side by side, and count the bars.
  \item Write \textbf{\ja{えき}} from memory and say \emph{eki}.
\end{enumerate}

\begin{tcolorbox}[breakable,colback=teal!4,colframe=teal!35!black,title={Before you move on}]
Stop after one accurate form.

Source: \href{https://www.unicode.org/charts/PDF/U3040.pdf}{Unicode Hiragana chart}; stroke count and order per \href{https://kanjivg.tagaini.net/kanjivg/kanji/0304d.html}{KanjiVG}.
\end{tcolorbox}
\section[\texorpdfstring{\ja{きのう} (kinō)}{きのう (kinō)}]{\ja{きのう} (kinō) — yesterday}

\label{lesson:JA-C133-kinou}



\begin{quote}
Write \textbf{\ja{き}} — \textbf{R1}, one lesson back. Then say \emph{a car} — \textbf{R2}, five lessons back.

\begin{itemize}
\raggedright
  \item \emph{Recall:} say \emph{a butcher} — \textbf{R4}, eighty lessons back
  \item \emph{Recall:} say \emph{tea} — \textbf{R3}, twenty lessons back
\end{itemize}
\end{quote}

\subsection*{You'll want to know: \ja{きのう}}
\textbf{\ja{きのう}} — \emph{kinō} — \textquotedblleft{}yesterday\textquotedblright{}.

The new sign is spent at once. \textbf{\ja{き}}, then \textbf{\ja{の}} and \textbf{\ja{う}}: three beats, \emph{ki–no–o}. The \textbf{\ja{う}} at the end stretches the \emph{o}, as it does in \textbf{\ja{ありがとう}}. You write every sign in it now.

\subsection*{Your turn}
\begin{itemize}
\raggedright
  \item \emph{Say it:} \emph{kinō}
  \item \emph{Say it:} \emph{eki}, then \emph{kinō}, and listen for the \emph{ki} in both
  \item \emph{Recall:} write \textbf{\ja{きのう}}, four strokes for the first sign
\end{itemize}

\begin{tcolorbox}[breakable,colback=teal!4,colframe=teal!35!black,title={Before you move on}]
What does \ja{きのう} mean? (\textquotedblleft{}Yesterday\textquotedblright{}.) Say it once more.
\end{tcolorbox}
\section[\texorpdfstring{\ja{いけ} (ike)}{いけ (ike)}]{\ja{いけ} (ike) — a pond}

\label{lesson:JA-C133-ike}



\begin{quote}
Say \emph{yesterday} — \textbf{R1}, one lesson back. Then write \textbf{\ja{る}} — \textbf{R2}, five lessons back.

\begin{itemize}
\raggedright
  \item \emph{Recall:} say \emph{hard} — \textbf{R4}, eighty lessons back
  \item \emph{Recall:} write \textbf{\ja{ゃ}}, small, in \textbf{\ja{おちゃ}} — \textbf{R3}, twenty lessons back
\end{itemize}
\end{quote}

\subsection*{You'll want to know: \ja{いけ}}
\textbf{\ja{いけ}} — \emph{ike} — \textquotedblleft{}a pond\textquotedblright{}.

Two beats: \emph{i–ke}. The first sign is the \textbf{\ja{い}} you already write. The second one is new, and the next lesson writes it.

\subsection*{Your turn}
\begin{itemize}
\raggedright
  \item \emph{Say it:} \emph{ike}
  \item \emph{Say it:} \emph{ike}, clapping two beats
  \item \emph{Recall:} point to the sign in \textbf{\ja{いけ}} you can already write, and the one you cannot
\end{itemize}

\begin{tcolorbox}[breakable,colback=teal!4,colframe=teal!35!black,title={Before you move on}]
What is \textquotedblleft{}a pond\textquotedblright{}? (\textbf{\ja{いけ}}, \emph{ike}.) How many beats? (\textbf{Two}.)
\end{tcolorbox}
\section[\texorpdfstring{\ja{け} (ke)}{け (ke)}]{\ja{け} — \ja{は}, without the loop}

\label{lesson:JA-W133-ke}



\begin{quote}
Four recalls, then the sign.

\begin{itemize}
\raggedright
  \item \emph{Recall:} say \emph{a pond} — \textbf{R1}, one lesson back
  \item \emph{Recall:} write \textbf{\ja{は}}: the left stroke, the bar, the loop
  \item \emph{Recall:} say \emph{correct} — \textbf{R4}, eighty lessons back
  \item \emph{Recall:} say \emph{a doctor} — \textbf{R3}, twenty lessons back
\end{itemize}
\end{quote}

\subsection*{Script — the shape}
\hlblockfigure{\detokenize{figures/JA-W133-ke-filmstrip.pdf}}{How \ja{け} is written, stroke by stroke}

\begin{quote}
\textbf{\ja{け}} — \emph{ke}
\end{quote}

Three strokes, two lifts:

\begin{enumerate}
\raggedright
  \item the left \textbf{vertical}, top to bottom; at the foot, turn and \textbf{flick} up to the right
  \item the \textbf{bar}, left to right
  \item the long \textbf{vertical} from the top, through the bar: run straight down, then \textbf{sweep} away to the lower left
\end{enumerate}

The first two strokes are the ones your hand already makes in \textbf{\ja{は}}. Only the last stroke differs:

\noindent
\begin{tabularx}{\linewidth}{@{}>{\raggedright\arraybackslash}X>{\raggedright\arraybackslash}X@{}}
\toprule
\textbf{you write} & \textbf{the last stroke ends} \\
\midrule
\textbf{\ja{は}}, \emph{ha} & in a small loop \\
\textbf{\ja{け}}, \emph{ke} & in a sweep to the lower left \\
\bottomrule
\end{tabularx}

With \textbf{\ja{い}} in front of it, you can now write the word from the last lesson: \textbf{\ja{いけ}}, \emph{ike}, a pond.

\subsection*{Your turn}
\begin{enumerate}
\raggedright
  \item Trace \textbf{\ja{け}} saying \textquotedblleft{}down and flick, bar, down and sweep\textquotedblright{}.
  \item Write \textbf{\ja{は}} and \textbf{\ja{け}} side by side; check that only the first one ties a loop.
  \item Write \textbf{\ja{いけ}} from memory and say \emph{ike}.
\end{enumerate}

\begin{tcolorbox}[breakable,colback=teal!4,colframe=teal!35!black,title={Before you move on}]
Stop after one accurate form.

Source: \href{https://www.unicode.org/charts/PDF/U3040.pdf}{Unicode Hiragana chart}; stroke count and order per \href{https://kanjivg.tagaini.net/kanjivg/kanji/03051.html}{KanjiVG}.
\end{tcolorbox}
\section[\texorpdfstring{\ja{けさ} (kesa)}{けさ (kesa)}]{\ja{けさ} (kesa) — this morning}

\label{lesson:JA-C133-kesa}



\begin{quote}
Write \textbf{\ja{け}} — \textbf{R1}, one lesson back. Then say \emph{a station} — \textbf{R2}, five lessons back.

\begin{itemize}
\raggedright
  \item \emph{Recall:} say \emph{free time} — \textbf{R4}, eighty lessons back
  \item \emph{Recall:} say \emph{a little} — \textbf{R3}, twenty lessons back
\end{itemize}
\end{quote}

\subsection*{You'll want to know: \ja{けさ}}
\textbf{\ja{けさ}} — \emph{kesa} — \textquotedblleft{}this morning\textquotedblright{}.

The new sign is spent at once. \textbf{\ja{け}} with \textbf{\ja{さ}} after it: two beats, \emph{ke–sa}, and both signs are ones you write. \textbf{\ja{きのう}} was yesterday; \textbf{\ja{けさ}} is earlier today.

\subsection*{Your turn}
\begin{itemize}
\raggedright
  \item \emph{Say it:} \emph{kesa}
  \item \emph{Say it:} \emph{kinō}, yesterday, then \emph{kesa}, this morning
  \item \emph{Recall:} write \textbf{\ja{けさ}}, three strokes for the first sign and three for the second
\end{itemize}

\begin{tcolorbox}[breakable,colback=teal!4,colframe=teal!35!black,title={Before you move on}]
What does \ja{けさ} mean? (\textquotedblleft{}This morning\textquotedblright{}.) And \ja{きのう}? (\textquotedblleft{}Yesterday\textquotedblright{}.)
\end{tcolorbox}
\section[\texorpdfstring{\ja{いぬ} (inu)}{いぬ (inu)}]{\ja{いぬ} (inu) — a dog}

\label{lesson:JA-C133-inu}



\begin{quote}
Say \emph{this morning} — \textbf{R1}, one lesson back. Then write \textbf{\ja{き}} — \textbf{R2}, five lessons back.

\begin{itemize}
\raggedright
  \item \emph{Recall:} say \emph{kind} — \textbf{R4}, eighty lessons back
  \item \emph{Recall:} write \textbf{\ja{ょ}}, small, in \textbf{\ja{ちょっと}} — \textbf{R3}, twenty lessons back
\end{itemize}
\end{quote}

\subsection*{You'll want to know: \ja{いぬ}}
\textbf{\ja{いぬ}} — \emph{inu} — \textquotedblleft{}a dog\textquotedblright{}.

Two beats: \emph{i–nu}. The first sign is the \textbf{\ja{い}} you already write. The second one is new. It looks like \textbf{\ja{め}} with a small loop tied at the end, and the next lesson writes it.

\subsection*{Your turn}
\begin{itemize}
\raggedright
  \item \emph{Say it:} \emph{inu}
  \item \emph{Say it:} \emph{inu}, clapping two beats
  \item \emph{Recall:} point to the sign in \textbf{\ja{いぬ}} you cannot write yet, and say which sign it looks like
\end{itemize}

\begin{tcolorbox}[breakable,colback=teal!4,colframe=teal!35!black,title={Before you move on}]
What is \textquotedblleft{}a dog\textquotedblright{}? (\textbf{\ja{いぬ}}, \emph{inu}.) How many beats? (\textbf{Two}.)
\end{tcolorbox}
\section[\texorpdfstring{\ja{ぬ} (nu)}{ぬ (nu)}]{\ja{ぬ} — \ja{め}, with a loop}

\label{lesson:JA-W133-nu}



\begin{quote}
Five recalls, then the sign.

\begin{itemize}
\raggedright
  \item \emph{Recall:} say \emph{a dog} — \textbf{R1}, one lesson back
  \item \emph{Recall:} say \emph{yesterday} — \textbf{R2}, five lessons back
  \item \emph{Recall:} write \textbf{\ja{め}}: the short stroke, then the long one that loops
  \item \emph{Recall:} say \emph{polite} — \textbf{R4}, eighty lessons back
  \item \emph{Recall:} say \emph{a library} — \textbf{R3}, twenty lessons back
\end{itemize}
\end{quote}

\subsection*{Script — the shape}
\hlblockfigure{\detokenize{figures/JA-W133-nu-filmstrip.pdf}}{How \ja{ぬ} is written, stroke by stroke}

\begin{quote}
\textbf{\ja{ぬ}} — \emph{nu}
\end{quote}

Two strokes, one lift:

\begin{enumerate}
\raggedright
  \item a short \textbf{diagonal}, from the upper left down to the right
  \item start at the top and cut \textbf{down to the left} across the first stroke; loop round the lower left, rise and run over the \textbf{arch}, come down the right side, then tie a small \textbf{loop} at the bottom and flick out to the right
\end{enumerate}

Up to the bottom of the right side, this is \textbf{\ja{め}}. Only the ending differs:

\noindent
\begin{tabularx}{\linewidth}{@{}>{\raggedright\arraybackslash}X>{\raggedright\arraybackslash}X@{}}
\toprule
\textbf{you write} & \textbf{it ends with} \\
\midrule
\textbf{\ja{め}}, \emph{me} & an open curve \\
\textbf{\ja{ぬ}}, \emph{nu} & a small closed loop and a flick \\
\bottomrule
\end{tabularx}

With \textbf{\ja{い}} in front of it, you can now write the word from the last lesson: \textbf{\ja{いぬ}}, \emph{inu}, a dog.

\subsection*{Your turn}
\begin{enumerate}
\raggedright
  \item Trace \textbf{\ja{ぬ}} saying \textquotedblleft{}short stroke; down, arch, loop\textquotedblright{}.
  \item Write \textbf{\ja{め}} and \textbf{\ja{ぬ}} side by side; check that only the second one closes a loop.
  \item Write \textbf{\ja{いぬ}} from memory and say \emph{inu}.
\end{enumerate}

\begin{tcolorbox}[breakable,colback=teal!4,colframe=teal!35!black,title={Before you move on}]
Stop after one accurate form.

Source: \href{https://www.unicode.org/charts/PDF/U3040.pdf}{Unicode Hiragana chart}; stroke count and order per \href{https://kanjivg.tagaini.net/kanjivg/kanji/0306c.html}{KanjiVG}.
\end{tcolorbox}
\section[\texorpdfstring{\ja{ぬの} (nuno)}{ぬの (nuno)}]{\ja{ぬの} (nuno) — cloth}

\label{lesson:JA-C133-nuno}



\begin{quote}
Write \textbf{\ja{ぬ}} — \textbf{R1}, one lesson back. Then say \emph{a pond} — \textbf{R2}, five lessons back.

\begin{itemize}
\raggedright
  \item \emph{Recall:} ask for tea politely — \textbf{R3}, twenty lessons back
\end{itemize}
\end{quote}

\subsection*{You'll want to know: \ja{ぬの}}
\textbf{\ja{ぬの}} — \emph{nuno} — \textquotedblleft{}cloth\textquotedblright{}, a piece of woven material.

The new sign is spent at once. \textbf{\ja{ぬ}} with \textbf{\ja{の}} after it: two beats, \emph{nu–no}, and both signs are ones you write. In \textbf{\ja{いぬ}} the new sign comes last; in \textbf{\ja{ぬの}} it comes first.

\subsection*{Your turn}
\begin{itemize}
\raggedright
  \item \emph{Say it:} \emph{nuno}
  \item \emph{Say it:} \emph{inu}, then \emph{nuno}, and listen for the \emph{nu} in both
  \item \emph{Recall:} write \textbf{\ja{ぬの}}, two strokes for the first sign and one for the second
\end{itemize}

\begin{tcolorbox}[breakable,colback=teal!4,colframe=teal!35!black,title={Before you move on}]
What does \ja{ぬの} mean? (\textquotedblleft{}Cloth\textquotedblright{}.) And \ja{いぬ}? (\textquotedblleft{}A dog\textquotedblright{}.)
\end{tcolorbox}
\section[\texorpdfstring{\ja{へや} (heya)}{へや (heya)}]{\ja{へや} (heya) — a room}

\label{lesson:JA-C133-heya}



\begin{quote}
Say \emph{cloth} — \textbf{R1}, one lesson back. Then write \textbf{\ja{け}} — \textbf{R2}, five lessons back.

\begin{itemize}
\raggedright
  \item \emph{Recall:} write \textbf{\ja{を}} — \textbf{R3}, twenty lessons back
\end{itemize}
\end{quote}

\subsection*{You'll want to know: \ja{へや}}
\textbf{\ja{へや}} — \emph{heya} — \textquotedblleft{}a room\textquotedblright{}.

Two beats: \emph{he–ya}. The second sign is the \textbf{\ja{や}} you already write. The first one is new: a single low peak, like a roof. When you wrote \textbf{\ja{ほ}}, the row \emph{ha, hi, fu, he, ho} had one sign missing. This is it, and the next lesson writes it.

\subsection*{Your turn}
\begin{itemize}
\raggedright
  \item \emph{Say it:} \emph{heya}
  \item \emph{Say it:} \emph{heya}, pointing at the room you are in
  \item \emph{Recall:} point to the sign in \textbf{\ja{へや}} you can already write, and the one you cannot
\end{itemize}

\begin{tcolorbox}[breakable,colback=teal!4,colframe=teal!35!black,title={Before you move on}]
What is \textquotedblleft{}a room\textquotedblright{}? (\textbf{\ja{へや}}, \emph{heya}.) How many beats? (\textbf{Two}.)
\end{tcolorbox}
\section[\texorpdfstring{\ja{へ} (he)}{へ (he)}]{\ja{へ} — one stroke, one peak}

\label{lesson:JA-W133-he}



\begin{quote}
Four recalls, then the sign.

\begin{itemize}
\raggedright
  \item \emph{Recall:} say \emph{a room} — \textbf{R1}, one lesson back
  \item \emph{Recall:} say \emph{this morning} — \textbf{R2}, five lessons back
  \item \emph{Recall:} write \textbf{\ja{ほ}}, and say the row it belongs to
  \item \emph{Recall:} say \emph{square} — \textbf{R4}, eighty lessons back
\end{itemize}
\end{quote}

\subsection*{Script — the shape}
\hlblockfigure{\detokenize{figures/JA-W133-he-filmstrip.pdf}}{How \ja{へ} is written, stroke by stroke}

\begin{quote}
\textbf{\ja{へ}} — \emph{he}
\end{quote}

One stroke, and the pen never lifts:

\begin{enumerate}
\raggedright
  \item start low on the left and \textbf{rise} to a low peak
  \item turn over the peak and run \textbf{down to the right}
\end{enumerate}

The peak sits left of the middle, so the way down is longer than the way up. Keep the peak low and the whole sign wide.

On its own after a place word, \textbf{\ja{へ}} is also a particle, and then it is read \emph{e}, the way \textbf{\ja{は}} is read \emph{wa} as the topic marker. That use comes later.

With \textbf{\ja{や}} after it, you can now write the word from the last lesson: \textbf{\ja{へや}}, \emph{heya}, a room.

\subsection*{Your turn}
\begin{enumerate}
\raggedright
  \item Trace \textbf{\ja{へ}} saying \textquotedblleft{}up, over, down\textquotedblright{}.
  \item Write the whole row \textbf{\ja{は} \ja{ひ} \ja{ふ} \ja{へ} \ja{ほ}}, and say it.
  \item Write \textbf{\ja{へや}} from memory and say \emph{heya}.
\end{enumerate}

\begin{tcolorbox}[breakable,colback=teal!4,colframe=teal!35!black,title={Before you move on}]
Stop after one accurate form.

Source: \href{https://www.unicode.org/charts/PDF/U3040.pdf}{Unicode Hiragana chart}; stroke count and order per \href{https://kanjivg.tagaini.net/kanjivg/kanji/03078.html}{KanjiVG}.
\end{tcolorbox}
\section[(dialogue)]{The last four — every basic sign, written}

\label{lesson:JA-R133-last-four}



\begin{quote}
Write \textbf{\ja{へ}} — \textbf{R1}, one lesson back. Then say \emph{a dog} — \textbf{R2}, five lessons back.

\begin{itemize}
\raggedright
  \item \emph{Recall:} say \emph{pink} — \textbf{R4}, eighty lessons back
  \item \emph{Recall:} say \emph{there}, near the listener — \textbf{R3}, twenty lessons back
\end{itemize}
\end{quote}

\subsection*{Your turn — read the seven words}
Read each one aloud and say what it means.

Read these aloud:

\begin{itemize}
\raggedright
  \item \textbf{\ja{えき}} — a station
  \item \textbf{\ja{きのう}} — yesterday
  \item \textbf{\ja{いけ}} — a pond
  \item \textbf{\ja{けさ}} — this morning
  \item \textbf{\ja{いぬ}} — a dog
  \item \textbf{\ja{ぬの}} — cloth
  \item \textbf{\ja{へや}} — a room
\end{itemize}

Each of these words holds one of the four signs this chapter wrote. They were the last four basic hiragana the book had not yet written.

\subsection*{Your turn — write}
\begin{enumerate}
\raggedright
  \item Write \textbf{\ja{き}}, \textbf{\ja{け}}, \textbf{\ja{ぬ}} and \textbf{\ja{へ}} once each.
  \item Write \textbf{\ja{きのう}} and \textbf{\ja{けさ}} from memory, and say which one is earlier.
  \item Write \textbf{\ja{へや}}, and say \emph{heya}.
\end{enumerate}

\begin{tcolorbox}[breakable,colback=teal!4,colframe=teal!35!black,title={Before you move on}]
A station? (\textbf{\ja{えき}}, \emph{eki}.) Yesterday? (\textbf{\ja{きのう}}, \emph{kinō}.) A pond? (\textbf{\ja{いけ}}, \emph{ike}.) This morning? (\textbf{\ja{けさ}}, \emph{kesa}.) A dog? (\textbf{\ja{いぬ}}, \emph{inu}.) Cloth? (\textbf{\ja{ぬの}}, \emph{nuno}.) A room? (\textbf{\ja{へや}}, \emph{heya}.)
\end{tcolorbox}
